
\begin{center}
\begin{figure}[H]
\centering
\begin{tikzpicture}[
  box/.style={rectangle, draw, minimum size=7mm, inner sep=0pt, font=\small},
  thick
]

\foreach \x in {0,...,9} {
  \pgfmathtruncatemacro{\i}{\x+1}
  \node[
    draw,
    minimum width=7mm,
    minimum height=7mm,
    inner sep=0pt,
    outer sep=0pt
  ] (m-1-\i) at (\x*7mm+3.5mm,3.5mm) {\x};
}

% Popisky nad buňkami
\foreach \x in {0,...,9} {
  \pgfmathtruncatemacro{\i}{\x+1}
  \node[above=2pt of m-1-\i] {\scriptsize \x};
}

% Hodnoty v přihrádkách pro h(k) = k mod 10
\node[below=12pt of m-1-3, box] (v1212) {1212};
\draw[->] (v1212.north) -- (m-1-3.south);

\node[below=12pt of m-1-6, box] (v935) {935};
\draw[->] (v935.north) -- (m-1-6.south);

\node[below=12pt of m-1-9, box] (v1948) {1948};
\draw[->] (v1948.north) -- (m-1-9.south);

\node[below=45pt of m-1-9, box] (v1918) {1918};
\draw[->] (v1918.north) -- (v1948.south);

\node[below=78pt of m-1-9, box] (v1968) {1968};
\draw[->] (v1968.north) -- (v1918.south);

\node[below=12pt of m-1-10, box] (v1989) {1989};
\draw[->] (v1989.north) -- (m-1-10.south);

\end{tikzpicture}
\caption{Ukázka hešovací tabulky s $m = 10$ přihrádkami a hešovací funkcí $h(k) = k \bmod 10$. Letopočty 1948, 1918 a 1968 vedou do stejné přihrádky 8 a tvoří kolizi. Zde vyřešeno pomocí řetízků}
\end{figure}
\end{center}